\section{Intro: Probability}

\begin{definition}

    A \textbf{discrete probability space} is a set of elementary events often denoted as $\Omega = \{
    \omega_1,
    \omega_2, \dots, \omega_n \}$.
    \\

    Every \textbf{elementary event} $\omega_i$ has a probability $P[\omega_i] \in [0,1]$ such that
    $$\sum_{\omega_i \in \Omega} P[\omega_i] = 1$$
    
\end{definition}

\begin{remark}

    A finite probability space, in which every elementary event has the same probability is called a
    \textbf{Laplace-Space}. Observe that for a Laplace-Space $\Omega$ and an arbitrary event $E$ we have:
    $$P[E] = \frac{\vert E \vert}{\vert \Omega \vert}$$
    
\end{remark}

\begin{definition}

    An \textbf{Event} $E \subseteq \Omega$ is a set of elementary events. The probability of an event is
    defined as:
    $$P[E] = \sum_{\omega_i \in E} P[\omega_i]$$
    The complementary event $\bar{E} = \Omega \setminus E$ has probability:
    $$P[\bar{E}] = 1 - P[E]$$
    
\end{definition}

\begin{theorem}

    For two events $A, B$ in a probability space $\Omega$ we have:
    \begin{itemize}
    \item $P[\emptyset] = 0$ and $P[\Omega] = 1$
    \item $P[A] \in [0,1]$
    \item $P[\bar{A}] = 1 - P[A]$
    \item If $A \subset B$ then $P[A] \le P[B]$
\end{itemize}

\end{theorem}

\begin{theorem}

    For events $A_1, A_2, ... A_n$ it holds
    $$P[A_1 \cup A_2 \cup \dots \cup A_n] \leq \sum_{i=1}^n P[A_i]$$
    with equality iff $A_1, A_2, ... A_n$ are pairwise disjoint.
    
\end{theorem}

\begin{theorem}

\textbf{Inclusion-Exclusion Principle: } For events $A_1, A_2, ... A_n$ with $n \geq 2$ it holds
    $$P[\bigcup_{i=1}^n A_i] = \sum_{i=1}^n (-1)^{i+1} P[\bigcap_{j=1}^i A_j]$$
    
\end{theorem}

    To proof this theorem we look at an arbitrary $\omega \in \bigcup_{i=1}^n A_i$. If $\omega$ is in exactly
    $k$ of the events $A_i$, then it is counted $N_k$ times where
    $$N_k = k - \binom{k}{2} + \binom{k}{3} - \dots + (-1)^{k+1} \binom{k}{k}$$
    We show that $N_k = 1$ for all $k \in \{1, 2, \dots, n\}$. For that we use the binomial theorem:
    $$(x-y)^k = \sum_{i=0}^k \binom{k}{i} x^{k-i} (-y)^i$$
    Now let $x=1$ and $y=1$:
    $$(1-1)^k = \sum_{i=0}^k \binom{k}{i} (-1)^i = \binom{k}{0} - \binom{k}{1} + \binom{k}{2} - \dots + (-1)^k
    \binom{k}{k} = 1- N_k$$
    Therfore:
    $$N_k = 1 - (1-1)^k = 1$$

    \subsection{Conditional Probability}

\begin{definition}

    Let $A$ and $B$ be events with $P[B] > 0$. The \textbf{conditional probability} of $A$ given $B$ is
    defined as:
    $$P[A \mid B] = \frac{P[A \cap B]}{P[B]}$$
    
\end{definition}

\begin{example}

    Assume a family with two children from which we know that one of them is a girl ($B$). What is the
    probability that the other child is also a girl ($A$)?
    $$P[A \mid B] = \frac{P[A \cap B]}{P[B]} = \frac{1/4}{3/4} = \frac{1}{3}$$
    
\end{example}

\begin{theorem}

\textbf{Multiplication-Theorem: } Let $A_1, A_2, .. A_n$ be events. If $P[A_1 \cap A_2 \cap \dots \cap
    A_{n-1}] > 0$ then it holds:
    $$P[A_1 \cap A_2 \cap \dots \cap A_n] = P[A_1] \cdot P[A_2 \mid A_1] \cdot P[A_3 \mid A_1 \cap A_2]
    \cdot \dots \cdot P[A_n \mid A_1 \cap A_2 \cap \dots \cap A_{n-1}]$$
    
\end{theorem}

    We can prove the correctness of this theorem by using the definition of conditional probability:
    $$P[A_1] \cdot P[A_2 \mid A_1] \cdot P[A_3 \mid A_1 \cap A_2] \dots P[A_n \mid A_1 \cap A_2 \cap \dots \cap
    A_{n-1}]$$
    $$ = P[A_1] \cdot \frac{P[A_1 \cap A_2]}{P[A_1]} \cdot \frac{P[A_1 \cap A_2 \cap A_3]}{P[A_1 \cap A_2]}
    \dots \frac{P[A_1 \cap A_2 \cap \dots \cap A_n]}{P[A_1 \cap A_2 \cap \dots \cap A_{n-1}]}$$
    $$ = P[A_1 \cap A_2 \cap \dots \cap A_n]$$

    \begin{example}

\textbf{Birthday Problem: } What is the probability of two people in a room of $m$ having the same
    birthday? We find the probability
    that none of them have the same birthday this is $\frac{364}{365} \cdot \frac{363}{365} \cdot \dots
    \cdot \frac{365-m+1}{365}$.
    
\end{example}

\begin{theorem}

\textbf{Theorem of total probability: }Let $A_1, ... A_n$ be pairwise disjoint events and let $B \subseteq
    A_1 \cup ... \cup A_n$ then it
    holds:
    $$P[B] = \sum_{i=1}^n P[B \mid A_i] P[A_i]$$
    
\end{theorem}

\begin{example}

\textbf{Goat Problem: } Assume three doors of which we know that one hides a sportscar and the other two a
    goat. You pick one
    arbitrary door, then one of the other two opens and you see a goat. Should you now switch to the other
    door
    that is left or should you stay with the one you are?

    Let $A_1$ be the event that you picked the door with the car and $A_2$ the event that you are infront of
    the door with a goat. Further let $B$ be the event that you win if you change doors. Now observe
    tracking-tight
    $$Pr[B \mid A_1] = 0$$
    $$Pr[B \mid A_2] = 1$$
    $$Pr[A_1] = 1/3$$
    $$Pr[A_2] = 2/3$$

    Using the theorem above we get:
    $$Pr[B] = Pr[B \mid A_1] Pr[A_1] + Pr[B \mid A_2] Pr[A_2] = 0 \cdot 1/3 + 1 \cdot 2/3 = 2/3$$
    
\end{example}

\begin{theorem}

\textbf{Bayes' Theorem: }
    Let $A_1, A_2 ... A_n$ be pariwise disjoint events and $B \subseteq A_1 \cup A_2 ... \cup A_n$ with
    $P[B] > 0$. Then:
    $$P[A_i \mid B] = \frac{P[B \mid A_i] P[A_i]}{\sum_{j=1}^n P[B \mid
    A_j] P[A_j]}$$
    
\end{theorem}

    To prove bayes theorem we use the theorem of total probability:
    $$P[A_i \mid B] = \frac{P[B \cap A_i]}{P[B]} = \frac{P[B \mid A_i] P[A_i]}{\sum_{j=1}^n P[B \mid
    A_j] P[A_j]}$$

    \begin{example}

\textbf{Desease Test: }We test a person for a desease that has a probability of $P[A] = 0.01$ to occur.
    Let $A$ be the event that the tested person has the desease. Further let $B$ be the event that the test
    was positive. Assume we know the sensitivity of the test $P[B \mid A] = 0.72$ and the specificity
    $P[\bar{B} \mid \bar{A}] = 0.98$. We want to find the probability that a person has the desease under
    the condition that the test is positive:
    $$P[A \mid B] = \frac{P[B \mid A] P[A]}{P[B \mid A] P[A] + P[B \mid
    \bar{A}] P[\bar{A}]} = \frac{0.72 \cdot 0.01}{0.72 \cdot 0.01 + 0.02 \cdot 0.99} = \frac{0.0072}{0.027}
    \approx 0.267$$
    Here we see that even if the test is quite accurate, the probability that a person has the desease
    is still quite low. This is because the probability of the desease is very low. If the probability of
    the desease was higher, the probability of the test being positive would be higher.
    
\end{example}

\subsection{Independence}

\begin{definition}

    Two events $A$ and $B$ are \textbf{independent} if $P[A \cap B] = P[A] P[B]$ or equivalently $P[A \mid B]
    = P[A]$.
    
\end{definition}

\begin{definition}

    Events $A_1, A_2, \dots, A_n$ are independent if for all subsets $I \subseteq \{1, 2, \dots, n\}$ it
    holds:
    $$P[\bigcap_{i \in I} A_i] = \prod_{i \in I} P[A_i]$$
    
\end{definition}

\begin{lemma}

    Let $A_1, A_2 \dots A_n$ be events and let $A_i^{0} \coloneqq \bar{A_i}$ and $A_i^{1} \coloneqq A_i$ for
    $i \in \{1, 2, \dots, n\}$.
    Then $A_1, A_2 \dots A_n$ are independent iff for all $(s_1, \dots s_n) \in \{0,1\}^{n}$ it holds that:
    $$P[\bigcap_{i=1}^n A_i^{s_i}] = \prod_{i=1}^n P[A_i^{s_i}]$$
    
\end{lemma}

\begin{lemma}

    Let $A, B$ and $C$ be independent events then so are $A \cap B$ and $C$ and also $A \cup B$ and $C$.
    
\end{lemma}

